\section{Requisitos funcionales}

La Tabla~\ref{tab:requisitos-funcionales} recoge los requisitos funcionales del MVP y el estado real de implementación al cierre del taller. Cada requisito se verificó con pruebas automatizadas o con el recorrido de extremo a extremo descrito en la Sección~\ref{sec:pruebas-parte2}.

\begin{table}[H]
	\caption{Requisitos funcionales de la solución}
	\label{tab:requisitos-funcionales}
	\centering
	\footnotesize
	\setlength{\tabcolsep}{3pt}
	\renewcommand{\arraystretch}{1.2}
	\begin{tabularx}{\textwidth}{@{}c >{\raggedright\arraybackslash}p{0.20\textwidth} Y c >{\raggedright\arraybackslash}p{0.11\textwidth}@{}}
		\toprule
		\textbf{Código} & \textbf{Requisito} & \textbf{Descripción} & \textbf{Prioridad} & \textbf{Estado} \\
		\midrule
		RF-01 & Registro de usuarios & Registrar y reconocer al usuario por su \texttt{telegram\_id} sin duplicar registros. & Alta & Sí \\
		RF-02 & Vinculación por invitación & Generar códigos aleatorios de un solo uso que crean una relación en estado pendiente. & Alta & Sí \\
		RF-03 & Autorización entre contactos & Permitir que el propietario autorice o rechace la consulta de su ubicación. & Alta & Sí \\
		RF-04 & Envío de ubicación & Recibir la posición solo cuando el usuario pulsa el botón nativo de Telegram. & Alta & Sí \\
		RF-05 & Almacenamiento & Guardar latitud, longitud, precisión horizontal y fecha/hora en UTC. & Alta & Sí \\
		RF-06 & Visualización en mapa & Mostrar la última posición propia y la de los contactos autorizados sobre Leaflet/OpenStreetMap. & Alta & Sí \\
		RF-07 & Historial & Consultar posiciones pasadas propias y de terceros con autorización vigente. & Alta & Sí \\
		RF-08 & Acceso web & Canjear un código \texttt{/web} de un solo uso por un JWT para el panel. & Alta & Sí \\
		RF-09 & Revocación & Revocar autorizaciones y detener la compartición con efecto inmediato. & Alta & Sí \\
		RF-10 & Estado del servicio & Exponer un \emph{health check} con el estado de la base de datos y del bot. & Media & Sí \\
		\bottomrule
	\end{tabularx}
\end{table}

\section{Requisitos técnicos}

La solución se ejecuta como tres servicios Docker: \texttt{postgres}, \texttt{backend} (API y bot en el mismo proceso) y \texttt{frontend} (panel estático servido por Nginx). La Tabla~\ref{tab:componentes-tecnicos} resume las tecnologías con las versiones fijadas en el repositorio.

\begin{table}[H]
	\caption{Componentes técnicos y tecnologías implementadas}
	\label{tab:componentes-tecnicos}
	\centering
	\footnotesize
	\setlength{\tabcolsep}{3pt}
	\renewcommand{\arraystretch}{1.2}
	\begin{tabularx}{\textwidth}{@{}>{\raggedright\arraybackslash}p{0.16\textwidth} Y >{\raggedright\arraybackslash}p{0.19\textwidth} c@{}}
		\toprule
		\textbf{Componente} & \textbf{Función} & \textbf{Tecnología y versión} & \textbf{Estado} \\
		\midrule
		Bot & Comandos, teclados y recepción de ubicación en modo \emph{polling}. & \texttt{python-\allowbreak telegram-bot} 22.8 & Sí \\
		API REST & Validación, autenticación y reglas de acceso. & FastAPI 0.162.4, Uvicorn 0.54.0 & Sí \\
		ORM & Modelos, relaciones y consultas. & SQLAlchemy 2.1.8 & Sí \\
		Base de datos & Persistencia de usuarios, permisos y ubicaciones. & PostgreSQL 16, \texttt{psycopg2} 2.9.13 & Sí \\
		Autenticación & Firma y validación de JWT del panel. & PyJWT 2.15.1 & Sí \\
		Configuración & Secretos y parámetros por entorno. & \texttt{pydantic-settings} 2.16.2 & Sí \\
		Panel web & Login, mapa, historial y autorizaciones. & React 19.3.1, Vite 8, React Router 7.18.5 & Sí \\
		Mapa & Teselas y marcadores sin claves de API. & React Leaflet 5.0.0, Leaflet 1.9.4, OpenStreetMap & Sí \\
		Servidor web & Estáticos y proxy \texttt{/api} hacia el backend. & Nginx & Sí \\
		Despliegue & Orquestación local reproducible. & Docker 29.4.3, Compose v5.1.4 & Sí \\
		Pruebas & Suites de API, permisos y bot. & \texttt{pytest} 9.2.0, \texttt{httpx} 0.28.1 & Sí \\
		\bottomrule
	\end{tabularx}
\end{table}

\subsection{Bot de Telegram}

El bot se ejecuta dentro del proceso de FastAPI: se inicializa en el \emph{lifespan} de la aplicación y usa \texttt{start\_polling()}, de modo que no requiere dominio público ni un contenedor adicional. La Tabla~\ref{tab:comandos-bot} lista los comandos disponibles.

\begin{table}[H]
	\caption{Comandos del bot de Telegram}
	\label{tab:comandos-bot}
	\centering
	\footnotesize
	\renewcommand{\arraystretch}{1.15}
	\begin{tabularx}{\textwidth}{@{}>{\raggedright\arraybackslash}p{0.22\textwidth} Y@{}}
		\toprule
		\textbf{Comando} & \textbf{Función} \\
		\midrule
		\texttt{/start}, \texttt{/menu} & Registro o bienvenida y menú principal con teclado personalizado. \\
		\texttt{/vincular} & Genera un código de invitación o explica el flujo de vinculación. \\
		\texttt{/vincular CÓDIGO} & Canjea una invitación recibida y crea la relación pendiente. \\
		\texttt{/compartir} & Muestra el botón nativo para enviar la ubicación actual. \\
		\texttt{/mis\_ubicaciones} & Resume las últimas posiciones propias y el estado de compartición. \\
		\texttt{/contactos} & Lista los contactos vinculados y el estado de los permisos. \\
		\texttt{/permisos} & Muestra las autorizaciones vigentes con opción de revocarlas. \\
		\texttt{/detener} & Detiene la compartición y oculta las posiciones a terceros. \\
		\texttt{/telefono} & Captura opcional del teléfono declarado por el usuario. \\
		\texttt{/web} & Emite un código de un solo uso para iniciar sesión en el panel. \\
		\texttt{/ayuda} & Instrucciones de uso y límites del sistema. \\
		\bottomrule
	\end{tabularx}
\end{table}

El teclado principal incluye \texttt{KeyboardButton(request\_location=True)}: la ubicación solo se transmite cuando el usuario pulsa ese botón y confirma el envío en el selector nativo de Telegram.

\subsection{API REST y panel web}

La Tabla~\ref{tab:endpoints-api} resume los endpoints implementados. Todas las rutas de datos exigen el encabezado \texttt{Authorization: Bearer <token>}; sin él la respuesta es \texttt{401} y, sin permiso vigente, \texttt{403}.

\begin{table}[H]
	\caption{Endpoints de la API REST}
	\label{tab:endpoints-api}
	\centering
	\footnotesize
	\setlength{\tabcolsep}{3pt}
	\renewcommand{\arraystretch}{1.1}
	\begin{tabularx}{\textwidth}{@{}c >{\raggedright\arraybackslash}p{0.30\textwidth} Y >{\raggedright\arraybackslash}p{0.13\textwidth}@{}}
		\toprule
		\textbf{Método} & \textbf{Ruta} & \textbf{Propósito} & \textbf{Autenticación} \\
		\midrule
		GET & \texttt{/api/health} & Estado del backend, la base de datos y el bot. & No \\
		POST & \texttt{/api/auth/token} & Canjea el código \texttt{/web} por un JWT. & Código de un solo uso \\
		GET & \texttt{/api/users/me} & Identidad del usuario autenticado. & JWT \\
		GET & \texttt{/api/users} & Contactos vinculados y estado de permisos. & JWT \\
		GET & \texttt{/api/authorizations} & Autorizaciones del usuario en ambos roles. & JWT \\
		POST & \texttt{/api/authorizations} & Crea o reactiva una autorización vinculada. & JWT \\
		DELETE & \texttt{/api/authorizations/\{id\}} & Revoca una autorización (solo el propietario). & JWT \\
		GET & \texttt{/api/locations/me} & Historial propio de ubicaciones. & JWT \\
		GET & \texttt{/api/locations/shared} & Últimas posiciones compartidas con el usuario. & JWT + doble consentimiento \\
		GET & \path{/api/locations/history/:id} & Historial de un tercero autorizado. & JWT + doble consentimiento \\
		GET & \path{/api/locations/latest/:id} & Última posición de un tercero autorizado. & JWT + doble consentimiento \\
		POST & \texttt{/api/invitations} & Genera un código de invitación. & JWT \\
		GET & \texttt{/api/invitations} & Invitaciones emitidas por el usuario. & JWT \\
		GET/PUT & \texttt{/api/config} & Consulta o cambia el interruptor de compartición. & JWT \\
		\bottomrule
	\end{tabularx}
\end{table}

El panel web consume la API a través del proxy \texttt{/api} de Nginx y nunca envía un identificador de usuario para decidir a quién puede consultar: la identidad se toma del token firmado. La Figura~\ref{fig:arquitectura}, en la Sección~\ref{sec:arquitectura-parte2}, describe el flujo entre componentes.

\subsection{Requisitos de seguridad y privacidad}

La Tabla~\ref{tab:requisitos-seguridad} relaciona los requisitos de seguridad con la implementación y la forma de verificación. Las pruebas citadas se ejecutaron sobre la pila real y se detallan en la Sección~\ref{sec:pruebas-parte2}.

\begin{table}[H]
	\caption{Requisitos de seguridad y privacidad de la solución}
	\label{tab:requisitos-seguridad}
	\centering
	\footnotesize
	\setlength{\tabcolsep}{3pt}
	\renewcommand{\arraystretch}{1.2}
	\begin{tabularx}{\textwidth}{@{}c >{\raggedright\arraybackslash}p{0.19\textwidth} Y >{\raggedright\arraybackslash}p{0.22\textwidth}@{}}
		\toprule
		\textbf{Código} & \textbf{Requisito} & \textbf{Implementación} & \textbf{Verificación} \\
		\midrule
		RS-01 & Consentimiento explícito & Botón \texttt{request\_location} y confirmación nativa de Telegram. & Pruebas del bot con un \texttt{telegram.Location} real. \\
		RS-02 & Doble consentimiento & Autorización activa y compartición habilitada; si falta una, la API responde \texttt{403}. & \texttt{test\_locations.py}: sin permiso, con permiso, tras revocar y con compartición detenida. \\
		RS-03 & Identidad web & JWT con \texttt{sub = telegram\_id}, emitido tras un código \texttt{/web} de un solo uso y cinco minutos. & Flujo completo código--token y casos de código inválido o caducado. \\
		RS-04 & Invitaciones seguras & Códigos de ocho caracteres con \texttt{secrets}, un solo uso y 15 minutos de vigencia. & Pruebas de unicidad, canje, reutilización y caducidad. \\
		RS-05 & Revocación inmediata & \texttt{DELETE} solo por el propietario; el acceso deja de funcionar al instante. & Revocación propia y \texttt{403} al intentar revocar la de otro. \\
		RS-06 & Configuración y secretos & Token del bot, contraseña de PostgreSQL y clave JWT por variables de entorno; \texttt{.env} excluido de Git. & Revisión de configuración y del historial del repositorio. \\
		RS-07 & Exposición de servicios & PostgreSQL sin puerto publicado; solo accesible desde la red interna de Compose. & \texttt{docker compose ps} y prueba de conexión externa fallida. \\
		RS-08 & CORS y errores & Orígenes permitidos explícitos y respuestas \texttt{500} genéricas sin trazas internas. & Revisión de middleware y pruebas de autenticación. \\
		RS-09 & Retención y borrado & Función \texttt{purge\_older\_than} y política de \texttt{LOCATION\_RETENTION\_DAYS}. & Prueba de purga por antigüedad y purga manual documentada. \\
		RS-10 & Minimización & Solo se guarda el teléfono si el usuario lo envía y la precisión únicamente cuando Telegram la informa. & Pruebas con precisión informada y con precisión nula. \\
		\bottomrule
	\end{tabularx}
\end{table}
